\subsection{Coherence criteria for the modules
\texorpdfstring{$\sh H^0_{X/Z}(\sh F)$}{H0(X/Z,F)}}
\label{subsection:IV.5.11}

\begin{proposition}[5.11.1]
\label{IV.5.11.1}
Let $X$ be a locally Noetherian prescheme, let $Z$ be a subset of
$X$ stable under specialization, and let $\sh F$ be a coherent
$\sh O_X$-module.  Denote by $(x_\alpha)$ the family of points of
\[
  \operatorname{Ass}(\sh F)\cap(X-Z),
\]
and, for every $\alpha$, let $Y_\alpha$ be the reduced closed
subprescheme of $X$ whose underlying space is
$\overline{\{x_\alpha\}}$, and put
$Z_\alpha=Z\cap\overline{\{x_\alpha\}}$.  The following two
conditions are equivalent:
\begin{enumerate}[{\rm a)}]
  \item $\sh H^0_{X/Z}(\sh F)$ is a coherent $\sh O_X$-module.
  \item For every $\alpha$,
  $\sh H^0_{Y_\alpha/Z_\alpha}(\sh O_{Y_\alpha})$ is a coherent
  $\sh O_{Y_\alpha}$-module.
\end{enumerate}
These conditions imply:
\begin{enumerate}
  \item[{\rm c)}] For every $\alpha$,
  $\codim(Z_\alpha,Y_\alpha)\geq2$.
\end{enumerate}
Moreover, conditions {\rm a)}, {\rm b)}, and {\rm c)} are
equivalent if, in addition, either of the following properties
holds:
\begin{enumerate}[{\rm(i)}]
  \item Every point of $X$ has an open neighbourhood isomorphic
  to a subscheme of a regular scheme; one also says in this case
  that $X$ is \emph{locally immersible in a regular scheme}.
  \item For every $\alpha$, $Y_\alpha$ is universally catenary
  \sref{IV.5.6.2}, and its normalization
  $\widetilde Y_\alpha$ \sref[II]{II.6.3.8} is finite over
  $Y_\alpha$.
\end{enumerate}
\end{proposition}

\noindent\emph{Proof.}\quad
All the properties under consideration are local on $X$.  We may
therefore restrict ourselves to the case in which
$X=\Spec(A)$ is affine, $A$ is Noetherian, and
$\sh F=\widetilde M$, where $M$ is a finite $A$-module.  For
every $\alpha$, let $h_\alpha:Y_\alpha\to X$ be the canonical
injection.  Then
\[
  (h_\alpha)_*(\sh O_{Y_\alpha})=\sh G_\alpha
\]
is the $\sh O_X$-module corresponding to the quotient
$A/\mathfrak j_{x_\alpha}$; by the definition of
$\operatorname{Ass}(\sh F)$, this quotient is isomorphic to an
$A$-submodule of $M$.  Since
$\sh H^0_{X/Z}(\sh G_\alpha)$ is a quasi-coherent
$\sh O_X$-submodule of $\sh H^0_{X/Z}(\sh F)$
\sref{IV.5.9.2}, condition {\rm a)} implies that it is coherent.
On the other hand, the definition \sref{IV.5.9.1.2} gives
\[
  \sh H^0_{X/Z}(\sh G_\alpha)
  \simeq
  (h_\alpha)_*
  \bigl(\sh H^0_{Y_\alpha/Z_\alpha}(\sh O_{Y_\alpha})\bigr).
\]
Thus {\rm a)} implies {\rm b)}.

To prove the converse, it is enough to find a finite filtration
$(\sh F_i)_{0\leq i\leq n}$ by coherent $\sh O_X$-modules, with
$\sh F_0=\sh F$ and $\sh F_n=0$, such that
\[
  \sh H^0_{X/Z}(\sh F_i/\sh F_{i+1})
\]
is coherent for every $i$.  Indeed, the assertion then follows by
descending induction from the exact sequences
\[
  0\longrightarrow\sh F_{i+1}\longrightarrow\sh F_i
  \longrightarrow\sh F_i/\sh F_{i+1},
\]
the left exactness of $\sh H^0_{X/Z}$ \sref{IV.5.9.2}, and
\sref[0I]{0.5.3.3} and \sref[I]{I.6.1.1}.  In view of
\sref{IV.3.2.8}, it is therefore enough to treat the case in which
$\sh F$ is irredundant, that is,
$\operatorname{Ass}(M)=\{\mathfrak p\}$ consists of one element.
We use the following lemma.

\begin{lemma}[5.11.1.1]
\label{IV.5.11.1.1}
Let $A$ be a Noetherian ring and let $M$ be a finite $A$-module
such that $\operatorname{Ass}(M)=\{\mathfrak p\}$.  There is a
finite filtration $(M_i)_{0\leq i\leq r}$ of $M$, with
$M_0=M$ and $M_r=0$, such that every quotient
$M_i/M_{i+1}$ is isomorphic to a submodule of $A/\mathfrak p$.
\end{lemma}

\begin{proof}
The canonical homomorphism
$M\to M_{\mathfrak p}=N$ is injective
(Bourbaki, \emph{Alg\`ebre commutative}, Chapter~IV,
\S1, no.~2, Proposition~6).  Put
\[
  B=A_{\mathfrak p},
  \qquad
  \mathfrak m=\mathfrak pA_{\mathfrak p}.
\]
Then $\mathfrak m$ is the maximal ideal of $B$ and
$\operatorname{Ass}(N)=\{\mathfrak m\}$
(\emph{loc.\ cit.}, Proposition~5).

\oldpage[IV-2]{123}

Since $N$ is a finite
$B$-module, there is an integer $s$ such that
$\mathfrak m^sN=0$.  The filtration by the powers
$\mathfrak m^jN$ may be refined to a finite filtration
$(N_i)$ of $N$ whose successive quotients are isomorphic to
\[
  B/\mathfrak m=\operatorname{Frac}(A/\mathfrak p).
\]
Intersecting it with $M$ gives $M_i=M\cap N_i$.  Each quotient
$M_i/M_{i+1}$ is then isomorphic to a finite
$(A/\mathfrak p)$-submodule of
$\operatorname{Frac}(A/\mathfrak p)$, and every such module is
isomorphic to a submodule of $A/\mathfrak p$.
\end{proof}

\begin{proof}[Completion of the proof of {\hyperlink{ega4.c04.prop.5-11-1}{Proposition~5.11.1}}]
The lemma and the preceding filtration argument reduce the proof,
under condition {\rm b)}, to showing that
$\sh H^0_{X/Z}(\sh G)$ is coherent when
\[
  \sh G=(A/\mathfrak p)^{\sim}
\]
and $\mathfrak p$ is an associated prime of $M$.  If
$\mathfrak p=\mathfrak j_y$ with $y\in Z$, the support of
$\sh G$ is a closed subset contained in $Z$, because $Z$ is
stable under specialization.  Hence
$\sh H^0_{X/Z}(\sh G)=0$ by \sref{IV.5.9.1.2}.  Otherwise,
$\mathfrak p=\mathfrak j_{x_\alpha}$ for some $\alpha$, and
\[
  \sh G=(h_\alpha)_*(\sh O_{Y_\alpha}).
\]
The displayed isomorphism above and condition {\rm b)} then show
that $\sh H^0_{X/Z}(\sh G)$ is coherent.  This proves the
equivalence of {\rm a)} and {\rm b)}.

Condition {\rm a)} implies {\rm c)} by
\sref{IV.5.10.10}[ii].  It remains to prove, under either
hypothesis {\rm(i)} or {\rm(ii)}, that {\rm c)} implies {\rm b)}.
If $X$ satisfies {\rm(i)}, each $Y_\alpha$ does as well.
Consequently, in view of \sref{IV.5.9.3.1}, the required
assertion follows from {\hyperlink{ega4.c04.cor.5-11-2}{Corollary~5.11.2}} below.
\end{proof}

\begin{corollary}[5.11.2]
\label{IV.5.11.2}
Let $A$ be a Noetherian domain satisfying either of the following
hypotheses:
\begin{enumerate}[{\rm(i)}]
  \item $A$ is a quotient of a regular Noetherian ring.
  \item $A$ is universally catenary, and its integral closure
  $A'$ is a finite $A$-algebra.
\end{enumerate}
Then
\[
  A^{(1)}
  =\bigcap_{\substack{\mathfrak p\in\Spec(A)\\
                      \operatorname{ht}(\mathfrak p)=1}}
    A_{\mathfrak p}
\]
is a finite $A$-algebra.
\end{corollary}

\begin{proof}
Suppose first that {\rm(i)} holds, and put $X=\Spec(A)$.  The set
$U$ of points at which $X$ satisfies $(\mathrm S_2)$ is open by
\sref{IV.6.11.2}.\footnote{The reader may verify that
{\hyperlink{ega4.c04.cor.5-11-2}{Corollary~5.11.2}} is not used in the proof of
\sref{IV.6.11.2}.}
Put $Z=X-U$.  Then $\codim(Z,X)\geq2$.  Indeed, if
$x\in X$ and $\dim(\sh O_x)\leq1$, then
$\dim(\sh O_{x'})\leq1$ for every generization $x'$ of $x$.
Since $\sh O_{x'}$ is a domain, it has positive depth whenever
its dimension is positive, and therefore $X$ satisfies
$(\mathrm S_2)$ at $x$.

Thus $Z\subset Z^{(2)}(X)$.  If $j:U\to X$ is the canonical
injection, \sref{IV.5.9.1.2} gives
\[
  \sh H^0_{X/Z^{(2)}(X)}(\sh O_X)
  \simeq
  j_*\bigl(
    \sh H^0_{U/Z^{(2)}(U)}(\sh O_U)
  \bigr).
\]
The prescheme $U$ satisfies $(\mathrm S_2)$, so
\sref{IV.5.10.14} identifies the module inside the parentheses
with $\sh O_U$.  Since $\codim(Z,X)\geq2$, Chapter~III,
\S9 shows that $j_*(\sh O_U)$ is coherent.  Finally,
\sref{IV.5.9.3.1} identifies the corresponding $A$-module with
$A^{(1)}$, proving the assertion under {\rm(i)}.

Now suppose that {\rm(ii)} holds.  The ring $A'$ is a Noetherian
normal domain, hence it is the intersection of its local rings
$A'_{\mathfrak p'}$ as $\mathfrak p'$ ranges over the
height-one prime ideals of $A'$
(Bourbaki, \emph{Alg\`ebre commutative}, Chapter~VII,
\S1, no.~6, Theorem~4).  For such a prime $\mathfrak p'$, put
\[
  \mathfrak p=\mathfrak p'\cap A,
  \qquad
  S=A-\mathfrak p.
\]

\oldpage[IV-2]{124}

The ring $A'_{\mathfrak p'}$ is the localization of the finite
$A_{\mathfrak p}$-algebra $S^{-1}A'$ at
$S^{-1}\mathfrak p'$.  The latter prime lies above the maximal
ideal $\mathfrak pA_{\mathfrak p}$ and is therefore maximal.
Moreover, $A_{\mathfrak p}$ is universally catenary by
\sref{IV.5.6.3}[i].  Formula \sref{IV.5.6.10} consequently gives
\[
  \dim(A_{\mathfrak p})
  =\dim(A'_{\mathfrak p'})
  =1.
\]
It follows that $A^{(1)}\subset A'$.  Since $A'$ is a finite
$A$-module and $A$ is Noetherian, its submodule $A^{(1)}$ is
finite as well.
\end{proof}
